\section{模型检验与评价}
\label{sec:model-validation}

\subsection{电量、费用与信息使用的检验}

模型求解后，另编程序根据实际运行记录重算供需平衡与储电量变化，不直接引用优化器报告的残差。对每个时段分别计算
\begin{equation}
\begin{aligned}
  r_t^{\rm bal}&=h_t+e_t+v_t+d_t-\ell_t-c_t-z_t,\\
  r_t^E&=E_{t+1}-E_t-\eta_cc_t+d_t/\eta_d.
\end{aligned}
\label{eq:residuals}
\end{equation}
同时检查电量上下限、充放电功率、互斥关系和跨日连续性，以及紧急购电是否仅补负荷。问题一两类残差均小于$10^{-8}$ kWh，问题三全部策略的电量核验残差不超过$2\times10^{-5}$ kWh；需要补购的时段由外网补足，未出现最终负荷缺供。

费用按有效计划、净调整量和实际电价逐段重算，再与日汇总、论文表格及结果工作簿核对。信息使用方面，保存预测所用历史和各版预报，将未来实测及未发布预报替换后重新计算，检查此前决策是否变化，详见附录\ref{app:validation-details}。这些检验确认的是既定时标、收费与反馈假设下的数值一致性，不能证明假设只有一种解释，也不能把近似规划的最优解视为原随机问题的最优策略。

\subsection{对照设置与统计结果的解释}

比较策略时，采用相同实际供需、初始储电量和对应收费规则，并同时列出期末电量与统一残值修正费用。总费用用于衡量总体支出，紧急电量反映计划缺口，日紧急费用CVaR95反映高损失日期的程度。三项指标分别讨论，不把某一项改善解释为所有方面都占优。

评价期为2025年2月至12月，一月用于积累历史及比较参数。对两策略同一天的费用作差，再按连续七天为块反复抽样，构造日均差的95\%区间，以保留部分短期相关性。区间仅描述本年度波动；由于方法开发曾查看全年诊断，二月至十二月不属于完全独立的测试集，结论仍需跨年度数据检验。

\subsection{模型的优点}

\textbf{各问模型能够逐步衔接。}四问共用电量平衡与储电量变化方程，随信息和调整条件变化扩展费用模型。去掉问题三的外部预报和调整，可以复现问题二；问题四改用固定价格曲线时也可复现前问，便于判断费用差异来自哪个新增条件。

\textbf{计划余量有可解释的费用依据。}从单时段多买、少买和调整电量的代价出发，得到0.8分位数及$[0.7,0.9]$调整区间，再将其作为多时段规划参数。这样能够说明为何预留电量以及何时修改计划，也明确了这些阈值适用的条件。

\textbf{求解结果有多方面证据支持。}第一问用费用下界检验最优性，后续各问通过相同条件下的逐项对照考察储能、预报、调整和电价的作用。所有账单由实际运行重算，对费用差不显著、补购风险增加及对照更便宜的结果均予保留。

\subsection{模型的不足与改进}

物理模型将微网作为整体、效率视为常数，未考虑线路容量、电池老化与控制延迟。若用于实际设备，需要补充网络和寿命约束，并用更细时段的数据检验快速反馈近似。

分位数曲线与调整区间简化了求解，但未完整描述连续多个时段共同出现的预测误差，也未充分考虑将来可以再次调整计划的收益。有限规划时长还可能遗漏远端储能价值。改进时应联合生成价格与需求情景，用真实补购损失替代区间内的近似收益，并在相同评价条件下比较。

本文结果还以正电价、取消部分退款、净额结算及余电允许弃用为前提。若收费规则或交付要求改变，应先重新明确现金流和可执行动作，再修改模型。完成这些检验后，才宜扩展到多储能设备、风光联合供能或需求响应情形。

\section{结论}

本文以电量平衡和储能状态变化为基础，分别建立四种信息与交易条件下的微网购电模型。确定条件下，联合优化购电与充放电的最低费用为33801.48元，比无储能方案节省29.66\%，并使光伏全部得到利用。结果说明，储能需要兼顾分时价差、未来负荷与富余光伏，不能只按当前电价决定动作。

当供需未知时，按历史误差预留计划电量可以减少高价补购，但延长规划时长未必继续节费：48小时方案改变了日末电量安排，与相同参数24小时方案的费用差仍不显著。允许日内调整后，相对相同预报下仅制定0时计划的策略节省\QThreeSavingPlanOnlyPercent\%，其中调整机会及实际状态反馈是主要来源；18时作用最大，6时最弱，最新光伏预报的额外节费尚不显著。

实时价格下，仅改变结算价格就使原决策费用增加约5.5\%，价格预测能收回其中一部分。日内有偿调整比全天计划不变节省\QFourThreeVsTwoPercent\%，调整及状态反馈占总节省的\QFourAdjustShareOfSaving\%。因此，进一步优化应着重减少高价时段的供能缺口，同时检查少数日期的高额补购；预测精度应与实际费用、补购风险一起评价。
